\documentclass[11pt]{article}

% ======================================================================
%  CbD analysis of S&P 500 (WRDS/CRSP) crisis periods.
%  Full mathematics + layman boxes. Compiles with pdflatex.
% ======================================================================
\usepackage[T1]{fontenc}
\usepackage[utf8]{inputenc}
\usepackage[a4paper,margin=0.86in]{geometry}
\usepackage{amsmath,amssymb,mathtools}
\usepackage{xcolor}
\usepackage[most]{tcolorbox}
\usepackage{enumitem}
\usepackage{booktabs}
\usepackage{array}
\usepackage{hyperref}
\hypersetup{colorlinks=true,linkcolor=blue!45!black,urlcolor=blue!45!black}

% ---------------------- callout boxes ---------------------------------
\definecolor{laymanblue}{RGB}{215,232,252}
\definecolor{laymanframe}{RGB}{60,120,200}
\definecolor{mathgreen}{RGB}{228,248,228}
\definecolor{mathgreenframe}{RGB}{40,140,40}
\definecolor{warnbg}{RGB}{255,244,215}
\definecolor{warnframe}{RGB}{190,130,0}

\newtcolorbox{intuition}[1][In plain terms]{
  enhanced, breakable, colback=laymanblue, colframe=laymanframe,
  fonttitle=\bfseries\small, coltitle=white, title={$\bigstar$\enspace #1},
  boxrule=1pt, arc=3pt, left=6pt, right=6pt, top=4pt, bottom=4pt,
  before skip=7pt, after skip=7pt}

\newtcolorbox{derivationbox}[1][Derivation]{
  enhanced, breakable, colback=mathgreen, colframe=mathgreenframe,
  fonttitle=\bfseries\small, coltitle=white, title={$\Sigma$\enspace #1},
  boxrule=1pt, arc=3pt, left=6pt, right=6pt, top=4pt, bottom=4pt,
  before skip=7pt, after skip=7pt}

\newtcolorbox{cautionbox}[1][Caveat]{
  enhanced, breakable, colback=warnbg, colframe=warnframe,
  fonttitle=\bfseries\small, coltitle=white, title={$\triangle$\enspace #1},
  boxrule=1pt, arc=3pt, left=6pt, right=6pt, top=4pt, bottom=4pt,
  before skip=7pt, after skip=7pt}

\newcommand{\sgn}{\operatorname{sgn}}
\newcommand{\ind}[1]{\mathbf{1}\!\left\{#1\right\}}
\newcommand{\E}[1]{\langle ab\rangle_{#1}}
\newcommand{\sodd}{s_{\mathrm{odd}}}

\title{\textbf{Contextuality-by-Default Analysis of S\&P 500 Crisis Periods}\\[4pt]
\large Design, Mathematics, and the Post-Selection Interpretation\\
\normalsize CRSP/WRDS daily data}
\author{\small Bell\,/\,stocks project --- analysis specification}
\date{\small\today}

\begin{document}
\maketitle

% ======================================================================
\section{Objective and framing}
% ======================================================================

This document specifies a complete Contextuality-by-Default (CbD) analysis of S\&P~500
constituent pairs over crisis and calm periods, using CRSP daily data from WRDS. The analysis
has three deliverables:
\begin{enumerate}[leftmargin=1.6em,itemsep=1pt]
  \item a working CbD pipeline layered on the existing four-$S$ (CHSH) machinery;
  \item a \emph{deflation result}: a demonstration that apparent Bell-type violations in
        equity data are reproducible by classical processes once the endogenous definition of
        the measurement contexts is accounted for; and
  \item a \emph{behavioral characterization}: how the dependence structure differs between
        crisis and calm regimes --- the aggregate System-1/System-2 contrast --- stated as
        classical regime-dependence, not contextuality.
\end{enumerate}

\begin{intuition}[What we are testing, and what we are not]
We are \emph{not} hunting for ``quantum stocks.'' The structural problem (\S\ref{sec:postsel})
is that in market data the measurement \emph{conditions} are computed from the same numbers as
the measurement \emph{outcomes} --- like labeling years ``drought years'' using the harvest
itself and then claiming drought explains the harvest. Because of this, a positive
contextuality score here cannot certify anything non-classical. The goals are therefore to show
this explicitly, quantify it, and extract the real (classical) crisis/calm structure --- while
the genuine contextuality question is reserved for decision experiments where conditions are
set externally.
\end{intuition}

% ======================================================================
\section{Data: CRSP via WRDS (build plan)}
% ======================================================================

This is a verified, build-ready specification. Two practical facts dominate it. First, the
historical S\&P~500 constituent list is \emph{no longer available through Compustat}: S\&P
withdrew the old \texttt{comp.idxcst\_his} table from WRDS, so index membership must now come
from CRSP. Second, CRSP migrated to its Flat File Format 2.0 (``CIZ'') in 2022; the legacy
(SIZ/FIZ) tables were frozen after the December~2024 data cut, and the CIZ \texttt{\_v2} tables
are the sole forward source. All table names below are the current CIZ ones.

\begin{cautionbox}[Two routes that will silently fail]
(1)~\texttt{comp.idxcst\_his} --- the old Compustat index-constituents table; S\&P pulled it
from the WRDS platform, so it no longer returns S\&P~500 membership. Use the CRSP list instead.
(2)~Legacy CRSP tables \emph{without} the \texttt{\_v2} suffix are frozen at the December~2024
cut and receive no new data; any sample reaching past 2024 must use the \texttt{\_v2} (CIZ)
tables. Both failure modes are quiet --- empty or stale results, not errors.
\end{cautionbox}

\textbf{Access.} Use the \texttt{wrds} Python package over the WRDS PostgreSQL connection; every
query below is a \texttt{db.raw\_sql(\dots)} call. Because CIZ renamed many columns, confirm
exact field names once with \texttt{db.describe\_table(...)} before relying on them.
\begin{verbatim}
    import wrds
    db = wrds.Connection(wrds_username="...")   # one-time .pgpass setup
    db.list_libraries()                         # confirm access
    db.describe_table("crsp", "dsf_v2")         # confirm CIZ column names
\end{verbatim}

\subsection{Point-in-time S\&P 500 membership}
Historical members of the index (SPX) are in \texttt{crsp.dsp500list\_v2} (daily; the
schema-qualified form is \texttt{crsp\_a\_indexes.dsp500list\_v2}). Each row is one membership
spell: \texttt{permno}, \texttt{mbrstartdt} (member start date), \texttt{mbrenddt} (member end
date). A security is a constituent on date $D$ iff
$\texttt{mbrstartdt}\le D\le\texttt{mbrenddt}$.
\begin{verbatim}
    spx = db.raw_sql("""
        select permno, mbrstartdt, mbrenddt
        from crsp.dsp500list_v2
    """)
\end{verbatim}
Applying this spell logic to \emph{each} window so the universe is the index composition
\emph{as of that date} is what removes survivorship and look-ahead bias --- using today's index
to study 2008 would silently condition on survival.

\begin{cautionbox}[Membership-date conventions (subtle, costly to miss)]
The end date of a \emph{currently active} member defaults to the last trading day of the latest
data cut and is revised forward at each annual update --- so current membership is
right-censored, not permanent. Because the list is daily, a membership test written against a
\emph{non-trading} day (e.g.\ Dec~31) can return nothing; anchor tests to trading days. Always
join on \texttt{permno}, never on ticker: tickers are reused and reassigned (Facebook$\to$Meta
is one of many).
\end{cautionbox}

\subsection{Daily returns and identifiers}
Daily returns are in \texttt{crsp.dsf\_v2} (the CIZ daily stock file), keyed by \texttt{permno}
and \texttt{date}. The holding-period return \texttt{ret} is adjusted for splits and
distributions, and --- a change from the legacy workflow --- in CIZ the \emph{delisting return is
already folded into the main return series}, so no separate delisting-file merge is required.
\begin{verbatim}
    ret = db.raw_sql("""
        select permno, date, ret, prc, vol, shrout
        from crsp.dsf_v2
        where date between '1990-01-01' and '2025-12-31'
          and permno in ( ... )       -- constituent PERMNOs for the window
    """)
\end{verbatim}
\textbf{Identifiers.} \texttt{permno} (permanent \emph{security} id) is the join key
everywhere; \texttt{permco} (permanent \emph{company} id) is used to avoid treating two share
classes of the same firm as an independent pair. \textbf{Universe filter.} The common-stock
screen (historically share codes 10/11 and exchange codes 1--3 = NYSE/AMEX/Nasdaq) now lives in
\texttt{crsp.stksecurityinfohist} under renamed CIZ fields (\texttt{sharetype},
\texttt{securitytype}, \texttt{primaryexch}, \dots); take the exact code values from the CRSP
CIZ cross-reference guide. For an S\&P~500-only study this screen is rarely binding (members are
common stock), but it matters if the universe is later broadened.

\subsection{Building the pair universe and data volume}
Per window: (i)~resolve constituent PERMNOs as of the window start via the spell logic;
(ii)~pull their daily \texttt{ret}; (iii)~form all unordered pairs. With $\sim$500 names a
window holds $\binom{500}{2}\approx1.25\times10^5$ pairs; across $\sim$420 monthly windows
(1990--present) that is $\sim$$5\times10^7$ pair-windows --- large but tractable. Cache the daily
return panel once (a few GB), compute each window vectorized over the return matrix, and persist
only the per-pair summaries ($\sodd,\Delta,\mathrm{CTX}$, the four cell counts), not the raw
cells. Membership churns by $\sim$20--30 names/year, so the universe is genuinely time-varying;
resolving it per window is mandatory, not a refinement.

\subsection{Continuous co-membership (pair eligibility)}
A pair contributes to a window only when \emph{both} names are constituents on that window's
trading days. Because membership turns over ($\sim$20--30 names/year, elevated in crises through
forced deletions), a naive construction would (i)~allow a pair's composition to change
\emph{within} a window when a name enters or exits mid-window --- contaminating the very
correlations being measured --- and (ii)~admit short-lived pairs whose mere availability is
correlated with the regime, a second-order echo of the endogeneity this analysis targets. We
therefore impose \textbf{continuous co-membership}: a pair $(A,B)$ is eligible for a window only
if both $A$ and $B$ are constituents for the \emph{entire} window (strict rule; a
$\ge\!95\%$-of-window-trading-days threshold is the permissive variant), with the co-membership
span recorded.

\begin{cautionbox}[Status of this rule: a standard principle, not a branded method]
There is no single citable ``continuous co-membership'' rule, and it would be an overstatement
to present one. What \emph{is} standard is its two ingredients, combined: (i)~a
\emph{traded-throughout eligibility screen} --- the formation/trading design of Gatev, Goetzmann
\& Rouwenhorst (2006), and the explicit requirement in later work that a security be present for
the whole formation window to be eligible (reviewed in Krauss 2017); and (ii)~\emph{point-in-time}
constituent reconstruction to remove survivorship and look-ahead bias. Our rule is exactly
(i)+(ii) applied to \emph{index membership} and to \emph{pairs}, rather than to single-stock
trading. Framing it this way is both defensible and accurate.
\end{cautionbox}

\subsection{Crisis dating (kept external to the pair)}
Window labels come from indicators \emph{outside} any pair's own returns:
(i)~\textbf{implied volatility} --- the CBOE VIX level (high/low by a fixed or rolling
threshold), available via WRDS (OptionMetrics) or directly from CBOE, and external to realized
constituent returns; (ii)~\textbf{NBER recession dates} --- macroeconomic, fully exogenous to
daily equity signs. The sample spans at least the 2000--02 dot-com unwind, the 2007--09
financial crisis, the 2011 euro-area stress, the 2015--16 drawdown, the 2020 COVID crash, and
the 2022 inflation/rates repricing.

\begin{cautionbox}[Why crisis dating must be external to the pair]
If the crisis label were computed from the same returns being analyzed, we would reproduce at
the regime level exactly the self-referential flaw this analysis exposes at the daily level. VIX
comes from S\&P index \emph{options} (a forward-looking, different instrument) and NBER dates
from macroeconomic aggregates --- both external to any constituent pair's own daily return path.
\end{cautionbox}

\subsection{Cleaning and quality}
(i)~Drop pair-days with a missing \texttt{ret} on either leg (halts, suspensions).
(ii)~Define $\sgn(0)$ explicitly: exact-zero daily returns occur for thinly traded names; treat
$\sgn(0)=0$ and exclude those days from the $\pm1$ series (or fix a convention) --- but be
consistent and report it. (iii)~Flag and down-weight ultra-thin names (many zero-volume days):
their signs are stale and inflate spurious agreement. (iv)~\texttt{ret} already incorporates
splits and dividends, so do not re-adjust prices; note that a negative \texttt{prc} in CRSP is a
quoted bid/ask midpoint (no closing trade), so use $|\texttt{prc}|$ if price is ever needed ---
though the analysis uses only \texttt{ret}, which is unaffected.

% ======================================================================
\section{The cyclic-4 system on a stock pair}
% ======================================================================

\textbf{Outcomes and contexts.} For a pair $(A,B)$ and day $t$, the outcomes are the return
signs $a_t=\sgn R_{A,t}$, $b_t=\sgn R_{B,t}\in\{-1,+1\}$. The context is the joint magnitude
regime $(x,y)$, with $x=\ind{|R_{A,t}|\ge\theta_A}$ mapped to $\{0,1\}$ ($0=$ large-move,
$1=$ small-move) and $y$ likewise for $B$. (The same construction runs with lagged-volatility
regimes; the magnitude version is the sharpest case of the endogeneity and is therefore the
right object for the deflation demonstration.)

\textbf{Contents and connections.} In CbD language the system has four \emph{contents}
(conceptual properties): $q_1=$ ``sign of $A$ when $A$ is in regime 0,'' $q_2=$ ``sign of $A$
in regime 1,'' $q_3=$ ``sign of $B$ in regime 0,'' $q_4=$ ``sign of $B$ in regime 1.'' Each
context $c_{xy}$ jointly measures one $A$-content and one $B$-content; each content appears in
exactly two contexts. This is a cyclic system of rank $n=4$: the same algebraic shape as the
CHSH experiment.

\begin{intuition}[Bunches and connections, on a real pair]
Take Deere (DE) and ADM --- two S\&P~500 names. Each trading day falls into one of four boxes,
by whether each stock had a big or small move that day:
\begin{center}\small
\renewcommand{\arraystretch}{1.25}
\begin{tabular}{l|cc}
 & \textbf{ADM big-move} & \textbf{ADM small-move}\\\hline
\textbf{DE big-move}   & $\E{00}$ & $\E{01}$\\
\textbf{DE small-move} & $\E{10}$ & $\E{11}$\\
\end{tabular}
\end{center}
In each box we record one number: over the window, how often do DE and ADM move the
\emph{same} direction on days that fall in that box. That single box --- two stocks measured
together under one condition --- is a \emph{bunch}. Now notice that ``DE's direction on its
big-move days'' appears in \emph{two} boxes (top row: ADM big, and ADM small). The comparison
of DE's average direction between those two boxes is a \emph{connection}; the four such
comparisons (two for DE, two for ADM) are exactly what $\Delta$ sums. Contextuality asks: can
one consistent story of how DE and ADM behave fill all four boxes at once --- allowing each
stock's average direction to drift a little between boxes (that drift is $\Delta$)? If yes, a
single classical description covers everything; if the four boxes cannot be reconciled even
after allowing the drift, the system is contextual.
\end{intuition}

\textbf{Empirical quantities.} Within each context $c_{xy}$, over a rolling window
$\mathcal{W}$ of width $w$ days:
\[
  \widehat{\E{xy}}
  =\frac{\sum_{t\in\mathcal{W}} a_t b_t\,\ind{(x_t,y_t)=(x,y)}}
        {N_{xy}},
  \qquad
  \widehat{\langle a\rangle}_{xy}
  =\frac{\sum_{t\in\mathcal{W}} a_t\,\ind{(x_t,y_t)=(x,y)}}{N_{xy}},
\]
with $N_{xy}$ the cell count, and $\widehat{\langle b\rangle}_{xy}$ analogously. Windows with
$\min_{xy}N_{xy}<N_{\min}$ (we use $N_{\min}=10$) are excluded: small cells mechanically
inflate the statistics.

% ======================================================================
\section{The CbD mathematics}
% ======================================================================

\subsection{The four CHSH combinations and \texorpdfstring{$\sodd$}{s\_odd}}

\[
  S_1=\E{00}+\E{01}+\E{10}-\E{11},\qquad
  S_2=\E{00}+\E{01}-\E{10}+\E{11},
\]
\[
  S_3=\E{00}-\E{01}+\E{10}+\E{11},\qquad
  S_4=-\E{00}+\E{01}+\E{10}+\E{11},
\]
\[
  \sodd \;=\; \max_{i\in\{1,\dots,4\}} |S_i|.
\]
Equivalently, $\sodd=\max\sum_{xy}\pm\E{xy}$ over sign patterns with an odd number of minus
signs --- hence the name.

\subsection{The inconsistency (direct-influence) term \texorpdfstring{$\Delta$}{Delta}}

Each content appears in two contexts; its two versions may have different means. The total
inconsistency is
\[
  \Delta
  =\big|\langle a\rangle_{00}-\langle a\rangle_{01}\big|
  +\big|\langle a\rangle_{10}-\langle a\rangle_{11}\big|
  +\big|\langle b\rangle_{00}-\langle b\rangle_{10}\big|
  +\big|\langle b\rangle_{01}-\langle b\rangle_{11}\big|.
\]

\begin{intuition}[What $\Delta$ measures]
$\Delta$ asks: does merely \emph{being in a different condition} shift a stock's average
up/down tendency? For example, does stock $A$ lean more negative on its large-move days than
its small-move days? (It does --- big down-days are bigger than big up-days.) This is ordinary
``the conditions differ'' behavior --- real, but shallow. CbD's point is that such shifts must
be measured and budgeted for \emph{before} any claim of deep structure.
\end{intuition}

\subsection{The Kujala--Dzhafarov--Larsson criterion}

For a cyclic system of rank $n$ with dichotomous outcomes, the system is noncontextual iff
\[
  \sodd \;\le\; (n-2) + \Delta \;\overset{n=4}{=}\; 2+\Delta,
  \qquad\text{i.e.}\qquad
  \boxed{\ \mathrm{CTX}\;\coloneqq\;\sodd-\Delta-2\;\le\;0.\ }
\]
$\mathrm{CTX}>0$ is the CbD definition of (true) contextuality. With consistent
connectedness ($\Delta=0$) this is exactly CHSH with its classical bound 2; quantum systems
reach $\sodd=2\sqrt2$; the algebraic maximum is 4.

\begin{derivationbox}[Why the classical bound is $2$ at $\Delta=0$]
Under a noncontextual (classical, common-description) model, all four contents have jointly
distributed representatives $A_0,A_1,B_0,B_1\in\{-1,+1\}$ and $\E{xy}=\langle A_xB_y\rangle$.
For any fixed realization, grouping the $S_1$ combination as
$A_0(B_0+B_1)+A_1(B_0-B_1)$: one bracket is $\pm2$, the other $0$, so each realization
contributes $\pm2$, and averaging gives $|S_1|\le2$; identically for $S_2,S_3,S_4$ (the single
minus sign only relocates which bracket is the difference). The KDL result extends this to
inconsistently connected systems: the achievable maximum of $\sodd$ under the best classical
stitching rises by exactly the marginal slack, i.e.\ by $\Delta$.
\end{derivationbox}

\subsection{A number you can check by hand}

\begin{derivationbox}[Watch the deflation happen]
Suppose one window yields these four co-movement numbers and eight conditional averages
(columns are the four regimes $00,01,10,11$):
\begin{center}\small
\renewcommand{\arraystretch}{1.2}
\begin{tabular}{l|cccc}
 & $00$ & $01$ & $10$ & $11$\\\hline
$\E{\cdot}$ \;(co-movement) & $+0.6$ & $+0.5$ & $+0.5$ & $-0.7$\\
$\langle a\rangle$ \;(DE avg.\ direction) & $-0.20$ & $-0.05$ & $+0.05$ & $+0.10$\\
$\langle b\rangle$ \;(ADM avg.\ direction) & $-0.15$ & $\phantom{+}0.00$ & $+0.05$ & $+0.10$\\
\end{tabular}
\end{center}
The largest combination is $S_1=0.6+0.5+0.5-(-0.7)=2.3$, so $\sodd=2.3>2$: a na\"ive
``violation.'' But the big-move cells lean negative (the leverage effect), so each stock's
average direction drifts across its two boxes:
\[
  \Delta = \underbrace{|{-}0.20-({-}0.05)|}_{0.15}
         + \underbrace{|0.05-0.10|}_{0.05}
         + \underbrace{|{-}0.15-0.05|}_{0.20}
         + \underbrace{|0.00-0.10|}_{0.10}
         = 0.50.
\]
The corrected score is $\mathrm{CTX}=\sodd-\Delta-2 = 2.3-0.50-2 = -0.20 \le 0$:
\emph{noncontextual}. The apparent violation was entirely the conditions shifting the averages
--- ordinary direct influence --- now subtracted out. On the post-selected market data of
\S\ref{sec:postsel} this deflation is the generic outcome.
\end{derivationbox}

% ======================================================================
\section{The post-selection no-go for market data}\label{sec:postsel}

The system above is structurally well-formed --- it \emph{is} a cyclic-4 system. The
obstruction to interpreting $\mathrm{CTX}>0$ as non-classicality is that the contexts are
defined by conditioning on $(|R_{A}|,|R_{B}|)$, which are functions of the very returns whose
signs are the outcomes. The context is therefore \emph{post-selected on an outcome-correlated
variable} --- the detection-loophole / measurement-dependence configuration. The classical
bound presumes contexts assigned independently of the system (``fair sampling''); that
presumption fails here by construction.

\begin{derivationbox}[A classical joint that reaches $\mathrm{CTX}>0$, indeed $\sodd=4$]
Construct a distribution for $(R_A,R_B)$ as follows. Let the magnitude pair
$(|R_A|,|R_B|)$ place positive probability on all four regime cells. Conditional on the cell,
draw the sign pair $(a,b)$ with marginals $\langle a\rangle_{xy}=\langle b\rangle_{xy}=0$ in
every cell (so $\Delta=0$ identically) and correlations
\[
  \E{00}=\E{01}=\E{10}=+c,\qquad \E{11}=-c,\qquad c\in(0,1].
\]
Unbiased $\pm1$ pairs with any prescribed correlation exist (mass
$\tfrac{1+c}{4}$ on equal signs, $\tfrac{1-c}{4}$ on opposite, etc.), so this is a valid,
purely classical joint distribution over $(R_A,R_B)$. Then
\[
  S_1=3c-(-c)=4c,\qquad \Delta=0,\qquad \mathrm{CTX}=4c-2>0\ \ \text{for } c>\tfrac12,
\]
reaching the algebraic maximum $\sodd=4$ at $c=1$ --- beyond the Tsirelson bound $2\sqrt2$,
which no genuine quantum system can exceed. A fully classical process did this. Hence
$\mathrm{CTX}>0$ on regime-sorted market data certifies nothing.
\end{derivationbox}

\begin{intuition}[Why $\Delta$ cannot rescue the test here]
$\Delta$ is a \emph{first-moment} correction: it budgets for the conditions shifting each
stock's \emph{average} tendency. But sorting days by move size distorts the
\emph{co-movement pattern} between the stocks at fixed averages --- a second-moment effect
that $\Delta$ never sees. In the construction above the averages are perfectly clean
($\Delta=0$) while the sorting drives the score to its arithmetic ceiling. The leak is in the
correlations, not the marginals, so the marginal correction cannot plug it.
\end{intuition}

\begin{cautionbox}[The two preconditions, and where they fail]
A CbD verdict reads as non-classicality only if (A) contexts are \emph{exogenous} ---
assigned independently of the measured system (measurement independence); and (B)
non-co-measurability arises from genuine \emph{disturbance}, not from \emph{selection}.
Observational equity data fails both: (A) the regimes are computed from the returns; (B) a
day sits in exactly one magnitude cell by sorting, not because measuring one content
physically precludes another. Designed decision experiments satisfy both, which is why true
contextuality has been demonstrated there (Cervantes--Dzhafarov 2018; Basieva et al.\ 2019)
and cannot be certified here.
\end{cautionbox}

% ======================================================================
\section{The analysis pipeline}
% ======================================================================

\textbf{Step 1 --- Estimation.} For every constituent pair and rolling window
($w=60$ trading days, stepped monthly; $\theta$ at the per-stock rolling median of $|R|$,
with fixed-quantile robustness sweeps): compute the four $\widehat{\E{xy}}$, the eight
conditional marginals, $\sodd$, $\Delta$, and $\mathrm{CTX}=\sodd-\Delta-2$. Record
$(N_{00},N_{01},N_{10},N_{11})$ and enforce $N_{\min}$.

\textbf{Step 2 --- Crisis/calm split.} Label each window crisis or calm by the external
indicator (\S2). Report, separately by regime: the cross-sectional distribution of $\sodd$,
$\Delta$, $\mathrm{CTX}$; the proportion of pairs with $\sodd>2$ (the na\"ive ``violation
rate''); and the proportion with $\mathrm{CTX}>0$ (the CbD-corrected rate).

\textbf{Step 3 --- The deflation demonstration (the central result).} Fit a classical
generator to each window: a two-state regime-switching model for magnitudes with a
sign-correlation that depends on the magnitude state (capturing crisis-coupling), plus an
asymmetric magnitude--sign link (capturing the leverage effect); no contextual ingredient
anywhere. Simulate, push through the identical pipeline, and compare the simulated and
empirical $(\sodd,\Delta,\mathrm{CTX})$ distributions. \emph{Reproduction} establishes the
artifact reading. Any empirical corner the classical generator cannot reach --- after the
controls below --- is flagged, not celebrated: the presumption remains artifact until the
generator class is shown inadequate.

\textbf{Step 4 --- Exogenous-context contrast.} Recompute with contexts defined by variables
\emph{not} derived from the pair's own returns: (i) the day's VIX level (high/low by rolling
median); (ii) scheduled macro-announcement days vs.\ ordinary days. Attenuation of
$\mathrm{CTX}$ under exogenous contexts is direct evidence that the endogenous sorting was the
source. (These contexts are shared, $x=y$, which empties the off-diagonal cells; the
comparison therefore uses the diagonal-cell statistics.)

\textbf{Step 5 --- Controls.} (i) Permutation placebo: shuffle regime labels within windows;
the $\mathrm{CTX}$ distribution must collapse below 0. (ii) Small-sample null: simulate
i.i.d.\ signs through the pipeline to quantify finite-$N$ inflation of $\sodd$.
(iii) Multiple comparisons: with $\sim$125{,}000 pairs per window, control the family-wise
rate (max-statistic permutation) before calling any single pair noteworthy.
(iv) Sector stratification: report Steps 2--3 within GICS sectors; concentrations are
candidates for the food-system follow-on, not contextuality claims.

\begin{intuition}[How to read the final table]
The headline exhibit is a two-row table --- crisis vs.\ calm --- with columns: na\"ive
violation rate ($\sodd>2$), CbD-corrected rate ($\mathrm{CTX}>0$), and the classical
simulator's matched rates. The expected pattern: the na\"ive rate is substantial and much
higher in crisis; the corrected rate is far smaller; and the simulator reproduces both. That
pattern, with the \S\ref{sec:postsel} construction beside it, \emph{is} the deflation result:
the crisis signal is real but classical --- stronger regime-dependent coupling --- and the
quantum-looking excess is the sorting artifact.
\end{intuition}

% ======================================================================
\section{Outputs and connection to the program}
% ======================================================================

Deliverables: (1) the pipeline (extending the four-$S$ codebase); (2) the deflation paper
exhibit set --- the table above, the $\mathrm{CTX}$ distributions, the classical-reproduction
overlay, and the exogenous-context contrast; (3) the crisis/calm behavioral characterization
($\Delta$, regime correlations, and their sector structure) as the aggregate
System-1/System-2 contrast, stated classically. The certifiable-contextuality question
proceeds in the decision/evacuation track, where contexts are experimenter-set and
incompatibility is genuine disturbance; this analysis is that track's methodological floor
and its motivating contrast.

\smallskip
\noindent\textbf{References} {\small(verified):} Kujala, Dzhafarov \& Larsson,
\emph{Phys.\ Rev.\ Lett.}\ \textbf{115}, 150401 (2015). \;Cervantes \& Dzhafarov,
\emph{Decision} \textbf{5}, 193--204 (2018). \;Basieva, Cervantes, Dzhafarov \& Khrennikov,
\emph{J.\ Exp.\ Psychol.: General} \textbf{148}, 1925--1937 (2019). \;Dzhafarov, Zhang \&
Kujala, \emph{Phil.\ Trans.\ R.\ Soc.\ A} \textbf{374}, 20150099 (2016).
\;\emph{Pair-eligibility methodology (confirm pagination):} Gatev, Goetzmann \& Rouwenhorst,
\emph{Rev.\ Financ.\ Stud.}\ \textbf{19}, 797--827 (2006); Krauss, \emph{J.\ Econ.\ Surveys}
\textbf{31}, 513--545 (2017).

\end{document}
